\documentclass{article} 
\usepackage{tabularray}
\UseTblrLibrary{booktabs}
\usepackage[a4paper]{geometry}
\usepackage{fancyhdr}
\pagestyle{fancy}
\lhead{Internationaler Terrorismus}
\rhead{März 2026}
\begin{document}
\section{Internationaler Terrorismus}
Gewalttätige Aktionen, die sich gegen den politischen oder gesellschaftlichen Status quo setzen, bei welchen auf unschuldige Dritte betroffen sein können, werden \emph{Terrorismus} genannt.
 
\subsection{Erscheinungsformen}
Diese werden in vier hauptsächliche Erscheinungsformen aufgeteilt: 
\begin{table}[htbp]
\begin{tblr}{
    colspec = {lXc}
  }
\toprule
 Form & Ziel & Beispiel \\
\midrule
 \emph{seperatistisch} & kämpfen für einen unabhängigen Staat einer religionen oder ethnischen Minderheit & IRA \\
 \emph{linksextremistisch} & kämpfen gegen wahrgenommene ungerechte und repressive wirtschaftliche Ordnungen & RAF \\
 \emph{rechtesextremistisch} & kämpfen gegen Minderheiten & NSU \\
 \emph{religiös} & kämpfen gegen säkulare Staaten und Gesellschaften, für eine theokratie & IS \\
\bottomrule%
\end{tblr}%
\end{table}%
\subsection{Bekämpfung}
Bei der Terrorismusbekämpfung mit Sicherheitsgesetzen gearbeitet werden, die den zuständigen Behörden mehr Befugnisse geben, welche eventuell von relevanz sein könnten. Dazu zählen Spionage durch Geheimdienste, Vorratsdatenspeicherung, Ausbildungen \& Rechte für Sicherheitsbehörden, usw. 
 
Des Weiteren kann \emph{Präventions}arbeit geleistet werden, die potentielle Terorristen frühzeitig deradikalisiert, \emph{Schutz}arbeit, um BürgerInnen und Infrastrukturen zu schützen, Terroristen verfolgt werden und eine angemessene \emph{Reaktion}, z.\,B. durch Krisenschutz, geplant werden.
\end{document}